% -------------------------------------------------------------------------
% Introdução
% -------------------------------------------------------------------------
\section{Introdução}
\label{sec:introducao}

As bolsas de valores funcionam como infraestruturas em que compradores e vendedores negociam ativos sob regras e supervisão, ajudando na formação de preços e na alocação de recursos \cite{mishkin2018,bodie2021}. No Brasil, a B3 --- Brasil, Bolsa, Balcão --- concentra negociação, registro, compensação e liquidação, sob marcos como a Lei do Mercado de Valores Mobiliários e a Lei das S.A. Os preços refletem expectativas sobre fluxos de caixa e risco, mas também exibem \textbf{volatilidade} e variação de \textbf{liquidez} que tornam a previsão de curto prazo um desafio tanto teórico quanto prático.

No horizonte \textbf{intradiário} aparecem questões clássicas de \emph{microestrutura}: periodicidades ao longo do dia, intervalos irregulares entre negócios e impacto do \textbf{spread} e do custo de imediatismo. Choques macroeconômicos e eventos externos podem deslocar as cotações rapidamente. Esse cenário motiva a busca por métodos preditivos que respeitem a ordem do tempo e as particularidades do mercado brasileiro.

Embora a ideia de que o mercado seja perfeitamente previsível seja indefensável, muitos investidores e pesquisadores procuram padrões que permitam melhorar decisões. Neste estudo, adotamos um modelo híbrido que combina \textbf{redes convolucionais} (CNN) e \textbf{redes recorrentes} do tipo \textbf{LSTM}: as primeiras capturam padrões locais em janelas deslizantes; as segundas, dependências de longo prazo. A escolha se apoia em evidências de que LSTMs representam melhor relações não lineares e sequências longas, enquanto convoluções extraem motivos de curto prazo; juntas, buscam-se maior robustez a ruído e a mudanças de regime.

Os alvos preditivos considerados são a \textbf{direção do retorno} na próxima barra (alta ou baixa), medidas de volatilidade intradiária e sinais relacionados a liquidez. Para assegurar comparabilidade e realismo, os resultados são confrontados com \textbf{modelos de referência} (\emph{baselines}) estatísticos (ARIMA, Prophet) e com validação \textbf{temporal deslizante} (\emph{walk-forward}), sem uso de informação futura no treino. O foco está em estabilidade e reprodutibilidade, e na pergunta: em que condições um modelo CNN+LSTM supera comparadores clássicos na previsão direcional intradiária (+1 barra, 15 min) em ações líquidas da B3, quando avaliado com métricas pós-custos?

A Figura~\ref{fig:delimitacao-tarefa} resume a tarefa: em uma série de preços intradiária, cada instante gera um vetor de atributos (preço, volume, indicadores técnicos) e o modelo estima a probabilidade de a próxima barra fechar acima ou abaixo da atual. O horizonte é de uma barra à frente, e as janelas de treino e teste avançam no tempo, refletindo o fluxo real do mercado.

\begin{figure}[htbp]
  \centering
  \includegraphics[width=0.75\linewidth]{delimitacao-tarefa.png}
  \caption{Previsão direcional intradiária: horizonte +1 barra (15 min). Fonte: autor.}
  \label{fig:delimitacao-tarefa}
\end{figure}
